% SEGMENT OLP-0628-S01
% Part: history
% Chapter: biographies
% Section: julia-robinson

\documentclass[../../../include/open-logic-section]{subfiles}

\begin{document}

\olfileid{his}{bio}{rob}

\olsection{ஜூலியா ராபின்சன்}

\olphoto{robinson-julia}{ஜூலியா ராபின்சன்}

ஜூலியா போமன் ராபின்சன் ஓர் அமெரிக்கக் கணிதவியலாளர். தீர்மானச்
சிக்கல்கள் பற்றிய பணிக்காகவும், குறிப்பாக ஹில்பர்ட்டின் பத்தாவது
சிக்கலுக்குத் தீர்வு காண்பதில் செய்த பங்களிப்புக்காகவும் அறியப்படுகிறார்.
1919 டிசம்பர் 8 அன்று மிசூரியின் செயின்ட் லூயிசில் பிறந்தார். சிறு
வயதிலேயே எண்கள் தமக்கு வியப்பூட்டியதாக ராபின்சன் நினைவுகூர்கிறார்
\citep[4]{Reid1986}. ஒன்பது வயதில் செந்நிறக் காய்ச்சல் ஏற்பட்டது;
அதைத் தொடர்ந்து வாதக் காய்ச்சல் மீண்டும் மீண்டும் வந்தது. இதனால்
நெடுங்காலம் படுக்கையிலேயே இருக்க வேண்டியதாயிற்று; கல்வியிலும்
பின்தங்கினார். தனிப்பயிற்சி ஆசிரியர்களின் உதவியால் பாடங்களில்
மீண்டும் முன்னேற முடிந்தாலும், நோயின் உடல்சார் விளைவுகள் அவருடைய
வாழ்க்கையில் நீடித்தன.

% SEGMENT OLP-0628-S02
சிறுவயதுச் சிரமங்களையும் மீறி, கணிதம், அறிவியல் ஆகியவற்றில் பல
விருதுகளுடன் மேல்நிலைப் பள்ளியை முடித்தார். சான் டியாகோ மாநிலக்
கல்லூரியில் பல்கலைக்கழகப் படிப்பைத் தொடங்கி, இறுதியாண்டு மாணவியாக
கலிபோர்னியா பல்கலைக்கழகத்தின் பெர்க்லி வளாகத்துக்கு மாறினார். அங்கே
கணிதவியலாளர் ரஃபேல் ராபின்சனின் தாக்கம் பெற்றார். இருவரும் நெருங்கிய
நண்பர்களாகி 1941-இல் திருமணம் செய்துகொண்டனர். ஆசிரியக் குழு
உறுப்பினரின் வாழ்க்கைத் துணை என்ற காரணத்தால், பெர்க்லியின் கணிதத்
துறையில் கற்பிக்க ராபின்சனுக்குத் தடை விதிக்கப்பட்டது. கணித
வகுப்புகளில் பதிவின்றிப் பங்கேற்பதைத் தொடர்ந்தாலும், பல்கலைக்கழகத்தை
விட்டு குடும்பம் தொடங்க விரும்பினார். ஆனால் திருமணமான சிறிது காலத்திலேயே
நிமோனியா ஏற்பட்டது. சிறுவயதில் ஏற்பட்ட வாதக் காய்ச்சலால் இதயத்தில்
கணிசமான தழும்புத் திசு உருவாகியிருப்பதாக அவரிடம் கூறப்பட்டது. அதன்
கடுமையைக் கருத்தில் கொண்டு, அவர் நாற்பது வயதுக்கு மேல் வாழமாட்டார்
என்று மருத்துவர் கணித்ததுடன், குழந்தைகள் பெற வேண்டாம் என்றும் அறிவுறுத்தினார்
\citep[13]{Reid1986}.

% SEGMENT OLP-0628-S03
நீண்ட காலம் மனச்சோர்வுடன் இருந்த ராபின்சன், இறுதியில் கணிதப் படிப்பைத்
தொடர முடிவு செய்தார். பெர்க்லிக்குத் திரும்பி, ஆல்ஃபிரட் தார்ஸ்கியின்
மேற்பார்வையில் 1948-இல் முனைவர் பட்டம் பெற்றார். மெய்யெண்களின் முதல்தரக்
கோட்பாடு தீர்மானிக்கத்தக்கது என்று தார்ஸ்கி காட்டியிருந்தார்; கோடலின்
பணியிலிருந்து இயல் எண்களின் முதல்தரக் கோட்பாடு தீர்மானிக்கவியலாதது
என்பது பின்தொடர்ந்தது. விகிதமுறு எண்களின் முதல்தரக் கோட்பாடு
தீர்மானிக்கத்தக்கதா என்பது முக்கியமான திறந்த சிக்கலாக இருந்தது.
தமது ஆய்வுரையில் அது தீர்மானிக்கவியலாதது என்று ராபின்சன் நிறுவினார்
\citeyearpar{Robinson1949}.

% SEGMENT OLP-0628-S04
தீர்மானச் சிக்கல்களில் ஆர்வம் கொண்ட ராபின்சன், அடுத்ததாக
ஹில்பர்ட்டின் பத்தாவது சிக்கலுக்குத் தீர்வு காண முயன்றார். 1900-இல்
டேவிட் ஹில்பர்ட் முன்வைத்த புகழ்பெற்ற 23 கணிதச் சிக்கல்களில் இதுவும்
ஒன்று. $3x^2 - 2y +3 = 0$ போன்ற, முழு எண் குணகங்களையுடைய பல்லுறுப்புச்
சமன்பாட்டுக்கு முழு எண்களில் தீர்வு உண்டா இல்லையா என்பதை முடிவுள்ள
நேரத்தில் விடையளிக்கும் படிமுறை உண்டா என்பதே பத்தாவது சிக்கல்.
இத்தகையவை \emph{டியோஃபான்டைன் சிக்கல்கள்} எனப்படும். தொடக்ககால
வெற்றிகளுக்குப் பின்னர், இதே சிக்கலில் பணியாற்றிய மார்ட்டின் டேவிஸ்,
ஹிலரி பட்னம் ஆகியோருடன் இணைந்தார். தெரியாத மாறிகள் அடுக்குகளாகவும்
வரக்கூடிய அடுக்குசார் டியோஃபான்டைன் சிக்கல்கள் தீர்மானிக்கவியலாதவை
என்று அவர்கள் காட்டினர். பின்னர் ``J.R.'' என்று அழைக்கப்பட்ட குறிப்பிட்ட
ஓர் ஊகம் உண்மையானால், ஹில்பர்ட்டின் பத்தாவது சிக்கல்
தீர்மானிக்கவியலாததாகும் என்றும் காட்டினர்
\citep{DavisPutnamRobinson1961}. ராபின்சன் 1960-கள் முழுவதும் இச்சிக்கலில்
பணியாற்றினார். 1970-இல் இளம் ரஷ்யக் கணிதவியலாளர் யூரி மதியாசெவிச்
இறுதியாக J.R. ஊகத்தை நிறுவினார். இணைந்த முடிவு இப்போது
மதியாசெவிச்--ராபின்சன்--டேவிஸ்--பட்னம் தேற்றம், சுருக்கமாக MRDP
தேற்றம், எனப்படுகிறது. மதியாசெவிச்சும் ராபின்சனும் நண்பர்களாகி பல
கட்டுரைகளில் கூட்டாகப் பணியாற்றினர். மதியாசெவிச்சுக்கு எழுதிய
கடிதத்தில் ராபின்சன், ``உண்மையில், ஆயிரக்கணக்கான மைல்கள் தொலைவில்
இருந்தபடியே நாம் சேர்ந்து பணியாற்றுவது, தனித்தனியாக யாராலும்
செய்யக்கூடியதைவிட அதிக முன்னேற்றத்தைத் தருகிறது என்பதில் எனக்கு
மிகுந்த மகிழ்ச்சி'' என்று ஒருமுறை எழுதினார் \citep[45]{Matijasevich1992}.

% SEGMENT OLP-0628-S05
அமெரிக்கக் கணிதச் சங்கத்தின் முதல் பெண் தலைவரும், தேசிய அறிவியல்
கழகத்துக்குத் தேர்ந்தெடுக்கப்பட்ட முதல் பெண்ணும் ராபின்சனே. இரத்தப்
புற்றுநோய் கண்டறியப்பட்ட பின்னர், 1985 ஜூலை 30 அன்று 65 வயதில்
இறந்தார்.

% SEGMENT OLP-0628-S06
\begin{reading}
ராபின்சனின் கணிதக் கட்டுரைகள் அவருடைய \textit{Collected
  Works} தொகுப்பில் கிடைக்கின்றன \citep{Robinson1996}; அதில்
தேசிய அறிவியல் கழகத்தின் வாழ்க்கை வரலாற்று நினைவுக்குறிப்பின்
மறுபதிப்பும் உள்ளது \citep{Feferman1994}. ராபின்சனின் அக்கா
கான்ஸ்டன்ஸ் ரீட், நேர்காணல்களின் அடிப்படையில் ``ஜூலியாவின்
சுயசரிதை'' ஒன்றை \citep{Reid1986} வெளியிட்டதுடன், முழு வாழ்க்கை
வரலாற்றையும் எழுதினார் \citep{Reid1996}. ராபின்சனையும்
ஹில்பர்ட்டின் பத்தாவது சிக்கலையும் பற்றிய குறும்பட ஆவணப்படத்தை
ஜார்ஜ் சிசெரி இயக்கினார் \citep{Csicsery2016}. மதியாசெவிச்,
ராபின்சனுடன் செய்த கூட்டுப் பணிகள், தமது பணியில் அவருடைய தாக்கம்
ஆகியவை பற்றிய சுருக்கமான நினைவுக்குறிப்புக்கு
\citep{Matijasevich1992}-ஐப் பார்க்கவும்.
\end{reading}

\end{document}
